\chapter{Data Sources, Acquisition, and Preprocessing}
\label{ch:data}

\section{Introduction}
\label{sec:data_intro}
The fidelity and reliability of any Earth observation machine learning pipeline depend fundamentally on the quality, calibration, and atmospheric correction of its input data. In satellite change detection, subtle variations in surface reflectance must be distinguished from sensor noise, orbital viewing geometry, and atmospheric scattering. This chapter details the remote sensing data sources utilized in TerraShift AI, the automated data acquisition and atmospheric quality masking pipeline, the integration of Google AlphaEarth 64-dimensional foundation embeddings, and the construction of the independently annotated ground-truth benchmark across four Saudi geographic corridors.

\section{Copernicus Sentinel-2 Multispectral Imagery}
\label{sec:data_sentinel2}
The primary optical remote sensing dataset utilized throughout this project is acquired by the European Space Agency's (ESA) Copernicus Sentinel-2 constellation, comprising two identical polar-orbiting satellites: Sentinel-2A (launched June 2015) and Sentinel-2B (launched March 2017).

\subsection{Sensor Characteristics and Spectral Bands}
The onboard Multispectral Instrument (MSI) captures sunlight reflected from the Earth's surface across thirteen spectral bands spanning the visible, near-infrared (NIR), and short-wave infrared (SWIR) electromagnetic spectrum (Table \ref{tab:sentinel2_bands}).

\begin{table}[htbp]
\centering
\caption{Copernicus Sentinel-2 MSI Spectral Band Specifications.}
\label{tab:sentinel2_bands}
\begin{tabularx}{\textwidth}{@{}lllXl@{}}
\toprule
\textbf{Band} & \textbf{Central $\lambda$ (nm)} & \textbf{Bandwidth (nm)} & \textbf{Spatial Res. (m)} & \textbf{Primary Physical Application} \\
\midrule
\textbf{B01} & 443 & 20 & 60 & Coastal aerosols, atmospheric correction \\
\textbf{B02} & 490 & 65 & 10 & Blue: Soil/vegetation differentiation \\
\textbf{B03} & 560 & 35 & 10 & Green: Green reflectance peak \\
\textbf{B04} & 665 & 30 & 10 & Red: Chlorophyll absorption \\
\textbf{B05} & 705 & 15 & 20 & Red Edge 1: Vegetation classification \\
\textbf{B06} & 740 & 15 & 20 & Red Edge 2: Canopy chlorophyll content \\
\textbf{B07} & 783 & 20 & 20 & Red Edge 3: Leaf area index (LAI) \\
\textbf{B08} & 842 & 115 & 10 & NIR: Biomass, water body boundary \\
\textbf{B8A} & 865 & 20 & 20 & Narrow NIR: Atmospheric water vapor \\
\textbf{B09} & 945 & 20 & 60 & Water vapor absorption \\
\textbf{B10} & 1375 & 30 & 60 & SWIR Cirrus: Sub-pixel cloud detection \\
\textbf{B11} & 1610 & 90 & 20 & SWIR-1: Soil moisture, built infrastructure \\
\textbf{B12} & 2190 & 180 & 20 & SWIR-2: Geology, carbonate detection \\
\bottomrule
\end{tabularx}
\end{table}

The constellation features a $290\text{ km}$ orbital swath width and a combined sun-synchronous revisit cycle of \textbf{5 days} at the equator, providing consistent multi-temporal observations under near-constant local solar illumination time ($\approx$ 10:30 AM local solar time).

\subsection{Level-1C vs. Level-2A Surface Reflectance}
A fundamental methodological requirement of this thesis is the exclusive use of \textbf{Level-2A (L2A) Bottom-Of-Atmosphere (BOA)} surface reflectance data:
\begin{itemize}
    \item \textbf{Level-1C (TOA):} Level-1C products represent Top-Of-Atmosphere reflectance, which includes the uncorrected effects of Rayleigh scattering, ozone absorption, water vapor attenuation, and aerosol backscatter. In desert environments, heavy seasonal aerosol loading (mineral dust haze) distorts L1C values, causing identical ground targets to exhibit temporal reflectance variances up to 30\%.
    \item \textbf{Level-2A (BOA):} Level-2A products are processed through the Sen2Cor atmospheric correction processor, deriving true physical surface reflectance [0, 1] after correcting for terrain elevation, gaseous absorption, and aerosol optical thickness. By evaluating L2A data, TerraShift eliminates spurious changes caused by temporal atmospheric shifts.
\end{itemize}

\section{Scene Classification Layer (SCL) Quality Masking}
\label{sec:data_scl}
Automated quality filtering is achieved via the Sentinel-2 Scene Classification Layer (SCL), generated by Sen2Cor at a 20-meter spatial resolution. SCL categorizes each pixel into one of twelve discrete radiometric classes (Table \ref{tab:scl_classes}).

\begin{table}[htbp]
\centering
\caption{Sentinel-2 Level-2A Scene Classification Layer (SCL) Pixel Typology.}
\label{tab:scl_classes}
\begin{tabular}{@{}lll@{}}
\toprule
\textbf{SCL Class Value} & \textbf{Class Label} & \textbf{Operational Action in TerraShift} \\
\midrule
0 & No Data / Defective & Invalid (Masked) \\
1 & Saturated or Defective & Invalid (Masked) \\
2 & Dark Area / Cast Shadows & Valid Surface \\
\textbf{3} & \textbf{Cloud Shadows} & \textbf{Invalid (Masked)} \\
4 & Vegetation & Valid Surface \\
5 & Bare Soils / Non-Vegetated & Valid Surface \\
6 & Water & Valid Surface \\
7 & Unclassified & Valid Surface \\
\textbf{8} & \textbf{Cloud Medium Probability} & \textbf{Invalid (Masked)} \\
\textbf{9} & \textbf{Cloud High Probability} & \textbf{Invalid (Masked)} \\
\textbf{10} & \textbf{Thin Cirrus} & \textbf{Invalid (Masked)} \\
11 & Snow / Ice & Valid Surface \\
\bottomrule
\end{tabular}
\end{table}

\subsection{Mathematical Formulation of the Clear-Land Mask}
For each pixel $(x, y)$, the invalid atmospheric mask $M_{\text{invalid}}(x, y)$ is defined as the logical union of cloud, shadow, and defective pixels:
\begin{equation}
    M_{\text{invalid}}(x, y) = \begin{cases} 
        1, & \text{if } \text{SCL}(x, y) \in \{0, 1, 3, 8, 9, 10\} \\ 
        0, & \text{otherwise} 
    \end{cases}
    \label{eq:scl_mask}
\end{equation}
For a bitemporal pair with acquisition epochs $T_1$ and $T_2$, the valid clear-land evaluation mask $V(x, y)$ requires valid ground observation at \textbf{both} dates:
\begin{equation}
    V(x, y) = \left(1 - M_{\text{invalid}}^{T_1}(x, y)\right) \land \left(1 - M_{\text{invalid}}^{T_2}(x, y)\right)
    \label{eq:bitemporal_validity}
\end{equation}
Any pixel where $V(x, y) = 0$ is strictly excluded from change inference and polygonization, preventing cloud edges or transient cloud shadows from spawning false change alerts.

\subsection{Scientific Disclosure on SCL Operational Limitations}
In the spirit of complete scientific transparency, we disclose that while SCL masking guarantees clear-land surface validity under standard operational criteria, it does not guarantee the complete absence of thin, diffuse sub-pixel cirrus or suspended mineral dust aerosols. In hyper-arid desert basins, dry dust haze can partially attenuate visible wavelengths without triggering SCL cloud thresholds (classes 8 or 9). Consequently, downstream algorithms cannot rely solely on optical reflectance; they must incorporate short-wave infrared indices and spatial size thresholds.

\section{Cloudless Annual Mosaic Composites (EOX)}
\label{sec:data_eox}
In addition to scene-level L2A acquisitions, the system integrates the EOX Sentinel-2 Cloudless global mosaic dataset. Produced by EOX IT Services GmbH, this dataset synthesizes billions of individual Sentinel-2 observations into seamless annual cloudless surface mosaics for each calendar year (2018 through 2024). The mosaic generation algorithm applies a temporal median compositing filter with statistical outlier rejection across the cloudless summer acquisition window, generating stable, radiometrically balanced true-color optical layers ideal for broad visual reconnaissance and regional context.

\section{Google AlphaEarth 64-D Latent Embeddings}
\label{sec:data_alphaearth}
To evaluate the frontier of self-supervised satellite representation learning, TerraShift AI incorporates Google AlphaEarth 64-dimensional annual satellite embeddings \citep{brown2025alphaearth}.

\subsection{Extraction via Google Earth Engine}
AlphaEarth embeddings are hosted on Google Earth Engine under the asset identifier \texttt{GOOGLE/SATELLITE\_EMBEDDING/V1/ANNUAL}. Each annual image composite contains 64 individual floating-point bands labeled $A_{00}$ through $A_{63}$ at an intrinsic $10\text{-meter}$ ground sampling distance.

During runtime, TerraShift interfaces with Earth Engine via the official Python \texttt{earthengine-api} client. Given an Area of Interest bounding box $\mathcal{B} = [\text{lon}_{\min}, \text{lat}_{\min}, \text{lon}_{\max}, \text{lat}_{\max}]$ and a target year $Y \in \{2021, 2024\}$, the system executes a spatial clipping and reduction query:
\begin{lstlisting}[language=Python, caption={Earth Engine AlphaEarth embedding query snippet.}]
import ee

ee_geom = ee.Geometry.BBox(west, south, east, north)
collection = (
    ee.ImageCollection("GOOGLE/SATELLITE_EMBEDDING/V1/ANNUAL")
    .filterDate(f"{year}-01-01", f"{year}-12-31")
    .filterBounds(ee_geom)
)
embedding_img = collection.first().select([f"A{i:02d}" for i in range(64)])
\end{lstlisting}

\subsection{Strict Prohibition of Synthetic Benchmark Data \& SHA-256 Verified Caching}
To guarantee uncompromising scientific integrity across all evaluations, \texttt{EarthEngineManager} enforces \texttt{allow\_synthetic=False} by default for all production and benchmark executions. When queried for any corridor, the system authenticates against Google Earth Engine (\texttt{project='terrashift'}) and verifies the 73-band tensor stack (64 AlphaEarth bands $A_{00}\dots A_{63}$ + 9 Dynamic World probability bands) against SHA-256 checksums stored in \texttt{services/ml/.cache/gee\_verified/}. If genuine Earth Engine data is unavailable, the pipeline immediately raises a \texttt{RuntimeError} rather than substituting synthetic embeddings. Deterministic fixture generation is isolated exclusively to explicitly flagged offline unit tests (\texttt{allow\_synthetic=True}).

\section{Google Dynamic World Land-Cover Probability Stacks}
\label{sec:data_dynamic_world}
To support semantic transition detection (both standalone \texttt{dynamic\_world\_lulc} and fused \texttt{alphaearth\_dynamic\_world}), the platform ingests Google Dynamic World V1 (\texttt{GOOGLE/DYNAMICWORLD/V1}) \citep{brown2022dynamicworld}. For any target year and bounding box, Dynamic World provides pixel-wise probability volumes across nine distinct land-use classes:
\begin{equation}
    \mathbf{P}_{\text{DW}}(x, y) = \left[ p_{\text{water}}, p_{\text{trees}}, p_{\text{grass}}, p_{\text{flood\_veg}}, p_{\text{crops}}, p_{\text{shrub}}, p_{\text{built}}, p_{\text{bare}}, p_{\text{snow}} \right]
    \label{eq:dw_probs}
\end{equation}
These probability tensors are extracted alongside AlphaEarth embeddings, providing explicit class transition deltas $\Delta p_{\text{built}} = p_{\text{built}}^{T_2} - p_{\text{built}}^{T_1}$ and $\Delta p_{\text{crops}} = p_{\text{crops}}^{T_2} - p_{\text{crops}}^{T_1}$.

\section{The Saudi Reference Benchmark Corridors}
\label{sec:data_saudi_benchmark}
A core scientific contribution of this thesis is the creation of the \textbf{Saudi Deployment Multi-Model Benchmark (v6.1)}. Spanning four geographic corridors across the Kingdom of Saudi Arabia, the benchmark encompasses \textbf{$91.19\text{ km}^2$} of diverse arid terrain, 19 parcel-level reference change events (coarse bounding-box envelopes curated from multi-date Sentinel-2 visual/spectral inspection and known Saudi landmarks), and four undisturbed negative control zones (Table \ref{tab:saudi_corridors_summary}).

\begin{table}[htbp]
\centering
\caption{Saudi Reference Evaluation Corridors Summary Specification.}
\label{tab:saudi_corridors_summary}
\resizebox{\textwidth}{!}{%
\begin{tabular}{@{}llllll@{}}
\toprule
\textbf{Corridor ID} & \textbf{Location / Region} & \textbf{Area ($\text{km}^2$)} & \textbf{Terrain / Land Cover} & \textbf{Change Typology} & \textbf{Evaluation Role} \\
\midrule
\texttt{riyadh\_sedra} & Riyadh (ROSHN Sedra) & 17.88 & Limestone desert regolith & Residential mass grading, paving & Core Development \\
\texttt{jeddah\_obhur} & Jeddah (North Obhur) & 18.31 & Coastal plain \& sabkha & Marine infill, road paving, boundary walls & Core Development \\
\texttt{aljouf\_sirhan} & Al-Jouf (Wadi As-Sirhan) & 26.82 & Hyper-arid agricultural basin & Circular pivot installation \& fallowing & Core Development \\
\texttt{alkharj\_logistics} & Al-Kharj Industrial Valley & 28.18 & Central desert \& oasis margin & Logistics warehouse pads, date palm parcels & \textbf{Held-Out Test Corridor} \\
\bottomrule
\end{tabular}%
}
\end{table}

\subsection{Corridor 1: Riyadh ROSHN Sedra ($17.88\text{ km}^2$)}
\begin{itemize}
    \item \textbf{Coordinates:} $[46.70^\circ\text{E}, 24.85^\circ\text{N}, 46.74^\circ\text{E}, 24.89^\circ\text{N}]$.
    \item \textbf{Environmental Context:} Located along the northern growth corridor of Riyadh, bordering natural limestone plateaus and dry tributary wadis.
    \item \textbf{Reference Events (5 discrete events):}
    \begin{enumerate}
        \item \textit{sedra\_mass\_grading\_sector1}: 52-hectare mass earthworks and residential sector excavation.
        \item \textit{sedra\_spine\_road\_paving}: Central dual-carriageway spine road grading and asphalt paving.
        \item \textit{sedra\_substation\_facility}: High-voltage electrical distribution substation foundation apron.
        \item \textit{sedra\_residential\_block\_a}: Phase 2 residential lot leveling and perimeter berming.
        \item \textit{sedra\_southern\_access\_junction}: Multi-lane highway access intersection and perimeter road.
    \end{enumerate}
    \item \textbf{Negative Control Zone:} Undisturbed natural limestone plateau and dry wadi ($[46.701, 24.875, 46.715, 24.889]$). Must generate \textbf{zero} false alerts.
\end{itemize}

\subsection{Corridor 2: Jeddah North Obhur Expansion ($18.31\text{ km}^2$)}
\begin{itemize}
    \item \textbf{Coordinates:} $[39.08^\circ\text{E}, 21.72^\circ\text{N}, 39.12^\circ\text{E}, 21.76^\circ\text{N}]$.
    \item \textbf{Environmental Context:} Western coastal plain north of Obhur Creek, characterized by coral-limestone substrate, coastal sabkhas, and marine infill.
    \item \textbf{Reference Events (4 discrete events):}
    \begin{enumerate}
        \item \textit{obhur\_coastal\_lot\_grading}: Large-scale residential subdivision lot leveling.
        \item \textit{obhur\_infill\_subdivision}: Residential infill plot boundary wall construction.
        \item \textit{obhur\_arterial\_junction}: Commercial boulevard grading and junction paving.
        \item \textit{obhur\_eastern\_access\_paving}: Dual carriageway access artery construction.
    \end{enumerate}
    \item \textbf{Negative Control Zone:} Undisturbed coastal sabkha and saline mudflat ($[39.082, 21.748, 39.098, 21.759]$), highly prone to false moisture-induced alarms in uncalibrated algorithms.
\end{itemize}

\subsection{Corridor 3: Al-Jouf Wadi As-Sirhan Basin ($26.82\text{ km}^2$)}
\begin{itemize}
    \item \textbf{Coordinates:} $[38.25^\circ\text{E}, 29.75^\circ\text{N}, 38.30^\circ\text{E}, 29.80^\circ\text{N}]$.
    \item \textbf{Environmental Context:} The agricultural heartland of northern Saudi Arabia, featuring commercial center-pivot circular irrigation fields bordered by the hyper-arid sands of the Great Nefud desert.
    \item \textbf{Reference Events (5 discrete events):} Five distinct 50-hectare circular center-pivot irrigation parcels (Alpha, Beta, Gamma, Delta, Epsilon) exhibiting either newly installed wheat/fodder cultivation or post-harvest fallowing.
    \item \textbf{Negative Control Zone:} Natural Great Nefud sand dunes ($[38.285, 29.775, 38.295, 29.795]$), testing algorithmic resilience against windblown sand transport.
\end{itemize}

\subsection{Corridor 4: Held-Out Al-Kharj Industrial Corridor ($28.18\text{ km}^2$)}
\begin{itemize}
    \item \textbf{Coordinates:} $[47.25^\circ\text{E}, 24.10^\circ\text{N}, 47.30^\circ\text{E}, 24.15^\circ\text{N}]$.
    \item \textbf{Methodological Role:} \textbf{Strict Held-Out Generalization Corridor}. All model architectures and hyperparameter thresholds ($\tau$, minimum polygon area) were completely frozen before evaluating Al-Kharj, providing an unbiased test of spatial generalization.
    \item \textbf{Reference Events (5 discrete events):}
    \begin{enumerate}
        \item \textit{alkharj\_logistics\_warehouse\_pad}: 35-hectare logistics warehouse foundation pad grading.
        \item \textit{alkharj\_industrial\_substation\_yard}: Industrial electrical substation and asphalt apron.
        \item \textit{alkharj\_commercial\_depot}: Open-air distribution depot and heavy vehicle storage yard.
        \item \textit{alkharj\_agricultural\_parcel\_surge}: Intensive date palm orchard and fodder parcel activation.
        \item \textit{alkharj\_industrial\_access\_strip}: Logistics access roadway corridor and drainage swale.
    \end{enumerate}
    \item \textbf{Negative Control Zone:} Undisturbed natural gravel desert plateau south of the industrial park ($[47.251, 24.101, 47.262, 24.108]$).
\end{itemize}

\section{Spatial Windowing and Co-Registration Protocol}
\label{sec:data_coregistration}
To ensure rigorous bitemporal alignment, all satellite raster tiles are projected into the Universal Transverse Mercator (UTM) coordinate system corresponding to their local geodetic zone (WGS84 UTM Zones 37N, 38N, and 39N). Multi-resolution bands (10m B02--B04, B08; 20m B11, SCL) are resampled to a unified 10-meter grid via bilinear interpolation for continuous reflectance and nearest-neighbor interpolation for categorical SCL masks. Sub-pixel spatial co-registration between epochs $T_1$ and $T_2$ is maintained by the ESA Copernicus Ground Segment, achieving root-mean-square spatial co-registration errors of less than $0.3\text{ pixels}$ ($<3\text{ meters}$), eliminating the risk of edge false alarms caused by misaligned grids.
\newpage
